% PENDING independent integration audit.  Written 2026-11-19 after checkpoint11.
% This file is not included in the paper until the parent promotes it.
\begin{saresult}{A.132}{Refuted}{The unrestricted real PL--TOP concordance comparison fails on the three-sphere}
Give the topological homeomorphism groups the compact-open topology,
and write $P_{\mathrm{TOP}}(S^3)$ for the singular complex of
\[
 \operatorname{Homeo}(S^3\times I\ \mathrm{rel}\ S^3\times\{0\}).
\]
The upper boundary is free. Write $P_{\mathrm{PL}}(S^3)$ for the
simplicial group of jointly finite PL concordances with the same
boundary convention. Then there is a weak equivalence of based spaces
\[
 P_{\mathrm{TOP}}(S^3)\simeq
 \Omega\bigl(\mathrm{TOP}(4)/\mathrm O(4)\bigr),
                                                        \tag{A.132.1}
\]
where $\mathrm{TOP}(4)=\operatorname{Homeo}(\mathbb R^4;0)$ is the
actual group of global homeomorphisms, and the quotient on the
right denotes $\operatorname{hofib}(B\mathrm O(4)\to
B\mathrm{TOP}(4))$. In particular,
\[
 \pi_2 P_{\mathrm{PL}}(S^3)=0,
 \qquad
 \pi_2 P_{\mathrm{TOP}}(S^3)=\mathbb Z/2.
                                                        \tag{A.132.2}
\]
Thus the natural real PL-to-TOP comparison is not a weak equivalence
in unrestricted dimension. Both groups have trivial $\pi_0$ and
$\pi_1$ in this example, so this discrepancy is genuinely a
higher-parameter discrepancy.
\end{saresult}
\begin{proof}
We prove the topological fibre model directly, including its relative
parameter convention. This avoids extending a PL germ theorem to TOP
without checking its hypotheses.

\smallskip\noindent
\textbf{The germ simplicial group and its two restriction maps.}
Let $D=D^4\subset\mathbb R^4$ be the closed unit ball, and put
\[
 A=\operatorname{Sing}\operatorname{Homeo}(D;0),\qquad
 T=\operatorname{Sing}\operatorname{Homeo}(\mathbb R^4;0).
\]
When $\mathrm O(4)$ or $\mathrm{SO}(4)$ occurs in a simplicial map
below, its singular complex is understood. Weak equivalences of
the resulting simplicial sets are equivalently weak equivalences
of their realizations.
Let $G$ be the following simplicial group of parameterized germs.
A $q$-simplex is the germ along $\{0\}\times\Delta^q$ of a
homeomorphism between open neighborhoods of that set in
$\mathbb R^4\times\Delta^q$, preserving the parameter projection
and fixing the zero section. Neighborhoods and openness are relative
to $\mathbb R^4\times\Delta^q$. Two representatives are equivalent
when they agree on a neighborhood of the whole zero section.
Composition, inverse, and the simplicial operators are defined on
germs. Compactness of $\Delta^q$ gives a common positive source
radius and a common positive image radius for each representative;
its inverse is jointly continuous on the latter neighborhood.
There are germ-restriction maps
\[
 r:A\longrightarrow G,\qquad g:T\longrightarrow G.
                                                        \tag{A.132.3}
\]

\smallskip\noindent
\textbf{A continuous local extension operator.}
We use the following precise consequence of
\cite[Theorem~5.1, p.~70; proof of Corollary~1.2, p.~79]
{EdwardsKirby1971}. Given concentric balls
$C\subset\operatorname{int}U\subset U\subset\operatorname{int}D$,
where $C$ and $U$ are compact, there is a neighborhood $\mathcal V$
of the inclusion in the space of embeddings $U\to\operatorname{int}D$
and a continuous map
\[
 \sigma:\mathcal V\longrightarrow
 \operatorname{Homeo}(D\ \mathrm{rel}\ \partial D)
\]
such that $\sigma(v)|_C=v|_C$ and $\sigma(\mathrm{id})=\mathrm{id}$.
The extensions may be chosen supported in a fixed compact subset of
$\operatorname{int}D$ after shrinking $\mathcal V$.
Indeed, Theorem~5.1 supplies a deformation, continuous in the input
embedding and fixed at the inclusion, to an embedding equal to the
identity on $C$ and unchanged near $\partial U$. Composing the input
with the inverse of this final embedding, and extending by the
identity, is exactly the local ambient extension in the proof on
p.~79. Its support stays inside the image of $U$; embeddings
sufficiently close to the inclusion have these images in a fixed
compact subset of $\operatorname{int}D$.
This is a continuous operator on an open set of embeddings, not
merely a separate extension selected for each parameter.

\smallskip\noindent
\textbf{Relative lifting over a horn, with uniform radii.}
We show that both maps in (A.132.3) are Kan fibrations. Let
$H=\Lambda_i^q\subset\Delta^q$, $q\geq1$, and suppose a horn
lift $a:H\to\operatorname{Homeo}(D;0)$ and a germ simplex $f$
are given with $r(a)=f|_H$. The proof for global homeomorphisms
is identical. Choose a representative $f_t$ on a common neighborhood
of $0$ for all $t\in\Delta^q$. Equality of germs on each of the
finitely many horn faces, followed by one common shrinking, gives
\[
 f_t=a_t\quad\hbox{on a fixed ball about $0$, for every }t\in H.
                                                        \tag{A.132.4}
\]
Choose a strong deformation $d_s:\Delta^q\to\Delta^q$, $0\leq s\leq1$,
with $d_0(\Delta^q)\subset H$, $d_1=\mathrm{id}$, and
$d_s|_H=\mathrm{id}$. Such a deformation exists because a simplex
collapses onto any horn. Set
\[
 a^0_t=a_{d_0(t)},\qquad
 k_{t,s}=f_{d_s(t)}\circ(a^0_t)^{-1}
\]
as germs. Joint continuity of $a^0$ and its inverse, compactness of
$\Delta^q\times I$, and (A.132.4) permit these germs to be
represented by one family of local homeomorphisms on a fixed ball,
with their inverses defined on another fixed ball. They satisfy,
on a common possibly smaller ball,
\[
 k_{t,0}=\mathrm{id}\quad(t\in\Delta^q),\qquad
 k_{t,s}=\mathrm{id}\quad(t\in H,\ s\in I).
                                                        \tag{A.132.5}
\]
Here the second equality is actual equality on the common ball:
$d_s(t)=d_0(t)=t$ on $H$ and (A.132.4) applies. It is not only
pointwise equality of unparameterized germs.

Choose a compact ball $U\subset\operatorname{int}D$ sufficiently
small that every $k_{t,s}^{-1}$ is defined on $U$ and its image
lies in the common source domain of every $k_{t,v}$.
This is possible by compactness and the joint inverse continuity
near the zero section. Shrink $U$ also so that (A.132.5) holds
on all relevant source and image neighborhoods. Choose a smaller
concentric compact ball $C\subset\operatorname{int}U$, and then
a still smaller ball $V$ such that
\[
 k_{t,s}(V)\subset\operatorname{int}C
 \quad\hbox{for all }(t,s)\in\Delta^q\times I.
                                                        \tag{A.132.6}
\]
The last choice follows from joint continuity and the fact that
every map fixes $0$.

On the fixed domain $U$, the increments
\[
 v_{t;s,s'}=k_{t,s}\circ k_{t,s'}^{-1}|_U
\]
depend continuously on $(t,s,s')$ in the embedding space wherever
they are used, and equal the inclusion when $s=s'$.
For the single neighborhood $\mathcal V$ in the local extension
operator above, compactness of the diagonal
$\{(t,s,s):t\in\Delta^q,s\in I\}$ gives $\eta>0$ such that
\[
 |s-s'|<\eta\quad\Longrightarrow\quad
 v_{t;s,s'}\in\mathcal V
 \quad\hbox{for all }t.
                                                        \tag{A.132.7}
\]
In particular these increments have images in $\operatorname{int}D$.
This supplies one time subdivision for all auxiliary parameters;
there is no parameter-dependent choice of subdivision or extension.

Choose $0=s_0<\cdots<s_N=1$ with mesh less than $\eta$.
Starting with $b_{t,0}=\mathrm{id}$, define for
$s_j\leq s\leq s_{j+1}$
\[
 b_{t,s}=
 \sigma(v_{t;s,s_j})\circ b_{t,s_j}.
                                                        \tag{A.132.8}
\]
The formulas agree at the subdivision points because
$\sigma(\mathrm{id})=\mathrm{id}$, so this is a continuous
family of disk homeomorphisms. By induction, (A.132.6) and the
extension property on $C$ give
\[
 b_{t,s}|_V=k_{t,s}|_V.
\]
Indeed, if $x\in V$, the input to the new extension is
$b_{t,s_j}(x)=k_{t,s_j}(x)\in C$.
The extensions fix $0$, since the increments do and their
extensions agree on $C$ containing $0$. For $t\in H$ every
increment in (A.132.8) is the identity by (A.132.5), hence
$b_{t,s}=\mathrm{id}$ for all $s$. Consequently
\[
 \widetilde a_t=b_{t,1}\circ a^0_t
\]
is a continuous lift of the given germ simplex and agrees exactly
with $a_t$ on $H$. Its germ is $f_t$ on a common ball, by one
further uniform shrinking under the family $a^0$ if necessary.
This proves the required horn lifting. For $T$, extend each
$b_{t,s}$ by the identity outside $D$ and use the same formula.
Thus arbitrary germ paths and compact auxiliary parameter families
are treated by finitely many small increments; no assumption that
the original germ family is close to the identity is made.

\smallskip\noindent
\textbf{The disk germ fibre has a free outer face.}
Let $K=\ker r$. For $0<\epsilon<1$, put
\[
 K_\epsilon=
 \operatorname{Sing}\operatorname{Homeo}(D\ \mathrm{rel}\ D_\epsilon),
 \qquad D_\epsilon=\epsilon D.
\]
Then $K$ is the filtered union of the $K_\epsilon$, because a
simplex in the identity germ fibre is the identity on a common
neighborhood of $0$. Restriction to the annulus identifies
$K_\epsilon$ with $P_{\mathrm{TOP}}(S^3)$: the inner sphere is
fixed pointwise, and the outer sphere is unrestricted. Extension
by the identity on $D_\epsilon$ is its continuous inverse.

For $0<\epsilon<\epsilon_0<1$, the inclusion
$K_{\epsilon_0}\to K_\epsilon$ is homotopic to an isomorphism.
Choose a radial disk homeomorphism $\rho$ taking
$D_{\epsilon_0}$ onto $D_\epsilon$ and a path $\rho_u$ from the
identity to $\rho$ for which
$\rho_u(D_{\epsilon_0})=D_{v(u)}$ with
$\epsilon\leq v(u)\leq\epsilon_0$. For example, take increasing
piecewise affine radius functions with this prescribed value at
$\epsilon_0$ and linearly interpolate the radius functions.
The homotopy $h\mapsto\rho_u h\rho_u^{-1}$ takes values in
$K_\epsilon$, starts at the inclusion, and ends at conjugation
by $\rho$, an isomorphism onto $K_\epsilon$.
Every finite simplicial sphere and every finite simplicial
nullhomotopy in $K$ uses finitely many simplices and therefore
one common radius. Filtered unions of these Kan complexes thus
give
\[
 K_{\epsilon_0}\xrightarrow{\simeq}K,
 \qquad K\simeq P_{\mathrm{TOP}}(S^3).
                                                        \tag{A.132.9}
\]
This identification imposes neither a fixed outer boundary nor
an additional factor recording a boundary parametrization.

\smallskip\noindent
\textbf{Global germs, with the component issue made explicit.}
The strict fibre $L=\ker g$ is weakly contractible. In fact, for
a compact parameter family $h_t$ in this fibre, all $h_t$ are the
identity on some $D_\epsilon$. The formula
\[
 c_s(h_t)(x)=
 \begin{cases}s^{-1}h_t(sx),&s>0,\\x,&s=0,\end{cases}
                                                        \tag{A.132.10}
\]
contracts the family to the identity in the compact-open topology.
On every fixed compact subset of $\mathbb R^4$, it is exactly the
identity for all sufficiently small $s$, uniformly in $t$; the
same assertion holds for inverses. For $0<s\leq1$ it is the
identity on $D_\epsilon$, and hence remains a family in $L$.
The formula fixes any parameters at which the input is the identity.

Let $T_0$ and $G_0$ be the identity components as simplicial sets.
The Kan fibration $g$ restricts to a Kan fibration
$T_0\to G_0$: a lift whose horn is in the identity component
remains there, because a simplex is connected. Every vertex of
$G_0$ can be reached by finitely many edges from the identity,
and the $q=1$ lifting just proved supplies a lift in $T_0$.
Together with (A.132.10), the fibre exact sequence proves
\[
 T_0\xrightarrow{\simeq}G_0.
                                                        \tag{A.132.11}
\]
No surjectivity assertion about all components of the germ group,
and no theorem globalizing an arbitrary isolated germ, is required.

\smallskip\noindent
\textbf{The disk group and the actual topological three-sphere group.}
Boundary restriction
$\operatorname{Homeo}(D;0)\to\operatorname{Homeo}(S^3)$ has the
radial-coning section. Its kernel is contracted by the relative
Alexander isotopy
\[
 a_u(x)=
 \begin{cases}u\,a(x/u),&\|x\|\leq u,\ u>0,\\
 x,&\|x\|\geq u\text{ or }u=0.
 \end{cases}
\]
These formulas agree on $\|x\|=u$ because $a$ fixes $\partial D$;
they fix $0$, and continuity at $u=0$ is uniform, with displacement
at most $2u$. The same bound holds for inverses. Splitting a
homeomorphism by its coned boundary restriction therefore proves
that boundary restriction is a homotopy equivalence.

The natural inclusion
$\mathrm O(4)\to\operatorname{Homeo}(S^3)$ is a homotopy
equivalence by \cite[Appendix, statement~(16), p.~606]{Hatcher1983}.
The original explicitly states the TOP version there; the
comparison discussion at statement~(3), p.~605 also explicitly
permits TOP in place of PL. Thus this is a statement about the
actual $C^0$ topological group, not an inference from the smooth
Smale theorem alone and not a statement about slice-wise PL maps.
It follows that
\[
 \mathrm O(4)\xrightarrow{\simeq}A,
 \qquad \mathrm{SO}(4)\xrightarrow{\simeq}A_+,
                                                        \tag{A.132.12}
\]
where $A_+$ is the orientation-preserving, hence identity,
component. An element of $K$ fixes a neighborhood of $0$ and is
orientation preserving, so the fibre of $r:A_+\to G_0$ is still
all of $K$.

Extend $a\in\operatorname{Homeo}(D;0)$ to a global homeomorphism
$e(a)$ by setting, outside $D$,
\[
 e(a)(r\omega)=r\,a|_{S^3}(\omega)\quad(r\geq1).
\]
This is a continuous group map, satisfies $g e=r$, and on
$\mathrm O(4)$ is the usual linear action on $\mathbb R^4$.
By (A.132.11), (A.132.12), and the fibre property already proved,
\[
 K\simeq\operatorname{hofib}(A_+\to T_0)
 \simeq\operatorname{hofib}(\mathrm{SO}(4)\to T_0)
 =\operatorname{hofib}_{1}(\mathrm O(4)\to T).
                                                        \tag{A.132.13}
\]
For the last equality, an orientation-reversing orthogonal map
cannot be joined to the identity in $T$. Every orientation-preserving
orthogonal map is in its identity component. Only those components
enter the homotopy fibre over $1$; no assertion about other
components of $T$ is needed.
For simplicial groups the contractible universal bundle identifies
$\Omega BG$ with $G$ up to weak equivalence. Applying this to the
homotopy fibre of $B\mathrm O(4)\to BT$ identifies the last space
in (A.132.13) with
$\Omega(\mathrm{TOP}(4)/\mathrm O(4))$. This and (A.132.9) prove
(A.132.1).

\smallskip\noindent
\textbf{The two homotopy calculations.}
The stabilization
\[
 \mathrm{TOP}(4)/\mathrm O(4)\longrightarrow\mathrm{TOP}/\mathrm O
\]
is $5$-connected by \cite[Theorem~8.7A, p.~128]{FreedmanQuinn1990}.
That source writes its comparison with $\mathrm{DIFF}(4)$;
for the precise global-group and orthogonal-group convention here,
we use the corrected modern account
\cite[\S1.1, Theorem~1.3 and the immediately following paragraph]
{GabaiGayHartmanKrushkalPowell2026} states the same
$5$-connected stabilization using exactly the global group
$\mathrm{TOP}(4)=\operatorname{Homeo}(\mathbb R^4;0)$ and the
classifying-space homotopy fibre convention used here.
It supplies the correction to Quinn's argument needed for the
full original theorem; the unchecked replacement criterion is
not used as an input.
By \cite[Essay~V, Theorem~5.5, p.~251]{KirbySiebenmann1977},
\[
 \pi_i(\mathrm{TOP}/\mathrm O)=
 \begin{cases}\mathbb Z/2,&i=3,\\0,&0\leq i\leq6,\ i\ne3.
 \end{cases}
\]
In particular (A.132.1) gives
$\pi_2P_{\mathrm{TOP}}(S^3)=\mathbb Z/2$ and vanishing
$\pi_0,\pi_1$.

For PL, set $X_j=\mathrm{PL}_j/\mathrm O_j$ as in A.128.
Its checked, jointly PL, free-top sphere fibre model is
\[
 P_{\mathrm{PL}}(S^3)\simeq
 \operatorname{hofib}(X_3\longrightarrow X_4).
\]
The low-degree formula
\cite[\S5, p.~37, formula~(3)]{BurgheleaLashof1974II} gives
$\pi_iX_j\cong\pi_i(\mathrm{PL}/\mathrm O)$ for $i<j$.
The latter groups vanish through degree six
\cite[Annex~C, p.~318]{KirbySiebenmann1977}.
The fibre exact sequence therefore has zero outer terms in
\[
 \pi_3X_4\longrightarrow\pi_2P_{\mathrm{PL}}(S^3)
 \longrightarrow\pi_2X_3,
\]
proving its middle group zero. The same argument in degrees
zero and one gives the asserted lower vanishing. The natural
PL-to-TOP inclusion consequently cannot induce an isomorphism
on $\pi_2$, proving the refutation.

For an explicit finite base, take the boundary of the abstract
$4$-simplex. The accompanying exact certificate checks its
$f$-vector $(5,10,10,5)$, every face link, integral boundary
Smith forms, and a primitive fundamental cycle: $53$ checks
pass. This computation certifies the finite PL three-sphere;
the existence and detection of the non-PL topological family
are proved by the fibre calculation and the cited stabilization
theorem, not asserted to have been numerically tested.
\end{proof}
